% Session 05 summaries
% Intended to be included from an existing main.tex.
% No document class, packages, or document environment are defined here.

\section{Allison -- Conceptual Models and the Cuban Missile Crisis}
\subsection{The problem of conceptual lenses}
Allison uses the Cuban Missile Crisis not only as an important historical case but also as a way to investigate how analysts explain governmental behavior. His central claim is methodological: explanations of foreign policy are never produced from facts alone. Analysts approach evidence through implicit \emph{conceptual models} that determine which facts appear important, what counts as a puzzle, which causal mechanisms deserve attention, and what kind of answer will count as an explanation. Different conceptual lenses can therefore generate different explanations of the same event without any change in the underlying historical evidence.

The article advances three general propositions. First, analysts of foreign and military policy normally rely on implicit conceptual models. Second, most conventional analysis is dominated by what Allison calls the \emph{Rational Policy Model} (Model I), in which a national government is treated as a unified actor selecting actions in pursuit of strategic objectives. Third, two alternative perspectives---the \emph{Organizational Process Model} (Model II) and the \emph{Bureaucratic Politics Model} (Model III)---can reveal determinants of governmental behavior that Model I obscures. The three perspectives are illustrated by applying each to the same central question: why did the United States respond to Soviet missiles in Cuba with a naval blockade rather than with another policy?

Allison's chess metaphor captures the difference among the models. Model I imagines a single chess player choosing moves strategically. Model II asks whether apparently coherent moves are instead generated by semi-independent organizations, each operating according to its own routines. Model III asks whether the moves are the result of bargaining among several players who share control over the pieces and pursue partly different objectives. The models thus direct attention respectively toward strategic choice, organizational routines, and political bargaining.

\subsection{Model I: Rational Policy}

Model I treats foreign-policy behavior as \emph{national choice}. The basic actor is a unitary national government with a coherent set of goals, a set of perceived alternatives, and estimates of the consequences of each alternative. A strategic problem creates threats and opportunities, and the government responds by selecting the action whose expected consequences best advance its objectives. The key elements are therefore goals and objectives, available options, expected consequences, and value-maximizing choice.

The model's dominant inference pattern runs backward from observed action to purpose. If a government performed action $X$, the analyst asks what objective would make $X$ a reasonable or optimal means. Explanation consists in reconstructing the rational connection between national objectives and the chosen action. Allison argues that this logic is widespread in international-relations scholarship even when it is not explicitly acknowledged. Different analysts may disagree about national objectives or constraints, but they frequently share the assumption that an important foreign-policy act is a calculated response by a unified state.

The model generates straightforward comparative propositions. If the costs of an option rise, or its probability of success falls, that option should become less likely; if its costs fall or its expected effectiveness rises, it should become more likely. Deterrence theory is a prominent example. A secure second-strike capability raises the expected costs of a nuclear first strike and should therefore reduce its attractiveness to a rational opponent.

Applied to the Cuban Missile Crisis, Model I depicts the blockade as the result of a systematic comparison of alternatives. The United States considered several broad courses of action. Doing nothing was rejected because the missiles were seen as militarily and politically significant and because failure to respond would undermine the credibility of American commitments. Pure diplomatic pressure appeared too slow and might force Washington into politically damaging concessions. A secret approach to Castro could not solve the problem because the missiles were controlled by the Soviet Union. A full invasion promised to remove both the missiles and Castro but risked direct combat with Soviet forces and a wider superpower confrontation. A surgical air strike offered speed but could not reliably destroy all missiles and risked appearing as a surprise attack. The blockade stood between passivity and immediate attack. It demonstrated resolve, exploited U.S. local superiority, slowed further Soviet deployments, and left Moscow time and room to retreat without immediate humiliation. From the Model I perspective, it was a value-maximizing form of limited escalation.

This first cut is powerful because it imposes coherence on the crisis. Yet precisely this coherence is what the other two models question. Governmental action may not be the product of a single calculation, and the options presented to leaders may themselves be shaped by organizational capabilities and political struggles.

\subsection{Model II: Organizational Process}

Model II opens the governmental ``black box.'' A government is not one actor but a conglomerate of large organizations. Leaders formally sit at the top, but most information, planning, preparation, and implementation are performed by organizations that possess established missions, routines, procedures, and repertoires. What Model I calls a national ``choice'' is therefore often better understood as an \emph{organizational output}.

Large organizations are indispensable because governments must monitor many problems and perform complex tasks simultaneously. They divide broad problems into narrower responsibilities and assign them to specialized units. This produces competence and regularity, but it also produces \emph{parochial priorities}. An organization tends to perceive problems through its own mission and capabilities. It develops standard operating procedures (SOPs) that allow ordinary members to handle recurring situations without reconsidering every problem from first principles. Complex actions are implemented through programs that coordinate many routines. Consequently, what leaders can actually order is constrained by the programs already available: governments can usually choose effectively only among actions for which organizations possess some existing capability and repertoire.

Model II therefore emphasizes continuity. The best clue to what an organization will do at time $t$ is often what it was doing at $t-1$. Organizational behavior changes incrementally because routines, budgets, professional norms, equipment, and division of labor cannot be redesigned instantly. Search is generally problem-directed and local: when performance falls below acceptable levels, organizations look first for solutions near existing routines. Learning and change occur, but they are usually gradual and shaped by existing organizational capacities.

This perspective changes the interpretation of the missile crisis. Soviet behavior need not be explained entirely by a single strategic calculation. Details of missile deployment, construction, camouflage, and military procedure may have reflected familiar Soviet organizational routines developed for other contexts and transferred to Cuba. Likewise, American intelligence did not simply ``decide'' when to discover the missiles. Detection depended on organizational procedures, reconnaissance routines, interpretive categories, and the availability of information through established channels.

The blockade itself also illustrates the autonomy of organizational implementation. Presidential leaders could decide on a quarantine in broad terms, but the Navy had to convert that political decision into an operational program. Naval doctrine and established procedures defined how ships would intercept, hail, inspect, and if necessary stop Soviet vessels. The resulting action was not infinitely adjustable to presidential wishes. Allison recounts the tension between Secretary of Defense Robert McNamara and Admiral George Anderson over how the blockade would actually be conducted. When McNamara demanded precise answers about interception procedures, Anderson relied on the Navy's established regulations and insisted that the Navy should run the blockade. The episode illustrates Model II's central point: even at the height of a crisis, implementation reflects organizational routines that political leaders cannot simply rewrite at will.

Model II thus qualifies the rational-choice story. Leaders may choose among broad alternatives, but the menu itself is organizationally generated, information is filtered through organizational channels, and execution is shaped by SOPs. Apparently puzzling details may therefore be explained not by grand strategic objectives but by the mundane fact that an organization did what it already knew how to do.

\subsection{Model III: Bureaucratic Politics}

Model III shifts from organizational routines to political bargaining. The leaders at the top of governmental organizations are not a single team with one objective. They are individual players occupying different positions, carrying different responsibilities, perceptions, interests, and commitments. Governmental action is therefore often a \emph{political resultant}: an outcome produced by ``pulling and hauling'' among players rather than a solution deliberately selected by a unified actor.

The model starts from three conditions. Power is shared among several participants; those participants disagree about what should be done; and the stakes are important enough that they fight to influence the decision. Position matters because responsibilities shape perceptions and priorities. This is summarized in the famous proposition that ``where you stand depends on where you sit.'' Military chiefs, diplomats, intelligence officials, and political advisers confront different organizational demands and therefore tend to define problems differently. The proposition also works vertically. Presidents, cabinet chiefs, senior staff, and lower-level officials face different constraints. Presidents must preserve flexibility and decide among competing claims; chiefs must build coalitions and secure presidential attention; lower-level officials struggle to place issues and proposals on the agenda.

Bureaucratic politics takes place through established \emph{action channels}: regular procedures that determine who participates in which decisions, at what stage, with what information and authority. A player's influence depends not only on formal position but also on bargaining advantages, expertise, control of information or implementation, access to the president, personal skill, willingness to use political resources, and the intensity of the player's preferences. Deadlines and previously made commitments matter because they narrow the space for bargaining and can allow decisions to emerge before every participant has reconsidered all alternatives.

The third cut at the Cuban Missile Crisis explains the blockade as the product of coalition formation inside the Kennedy administration. Important advisers initially favored different options. The Joint Chiefs and several senior figures strongly supported an air strike. McNamara increasingly regarded a blockade as a more controllable alternative. Robert Kennedy objected to a surprise air attack on moral and political grounds, fearing that it would make the United States resemble the perpetrators of Pearl Harbor. Theodore Sorensen also moved toward the blockade. The personal and political relationship of these advisers with President Kennedy mattered, as did the way different options were presented.

A decisive coalition gradually formed around McNamara, Robert Kennedy, and Sorensen. Meanwhile, advocates of an air strike included figures whose style and relationship with the president made their recommendations less persuasive to him. Information also entered the political game. The claim that an air strike could not be reliably ``surgical'' strengthened the blockade coalition, even though Allison notes that some of the information underlying this judgment was inaccurate or insufficiently examined. The final choice therefore cannot be reduced to a neutral ranking of fixed alternatives. It emerged from the timing of meetings, coalition building, personalities, institutional positions, arguments, and the distribution of influence among the players.

Model III also changes the meaning of governmental intention. An outcome can be coherent enough to count as government policy without having been intended in all its details by any one participant. Different actors may contribute different pieces of the final action for different reasons. The result can therefore be distinct from the preferred policy of every individual participant.

\subsection{Overall implications}

Allison does not argue that one model is universally correct and the others false. Rather, each highlights one class of determinants while holding others in the background. Model I emphasizes pressures and incentives in the international strategic environment. Models II and III examine the internal governmental machinery through which those pressures are perceived and converted into action. They are not mutually exclusive, and a fuller theory would need to specify which kinds of decisions are best explained by which mechanisms and how the mechanisms interact.

The broader lesson is a demand for analytical self-consciousness. Researchers should identify the assumptions embedded in their explanations rather than treating one conceptual lens as natural or neutral. The same event can look like a rational choice, an organizational output, or a political resultant depending on the model used. Models II and III are particularly useful as correctives to the tendency to infer large purposes from large events. They can improve explanation and prediction by directing attention to organizational capabilities, routines, intra-governmental disagreement, and bargaining.

At the same time, Allison recognizes limitations. Model I can rationalize almost any observed action if objectives are freely reconstructed. Model II risks explaining the present too heavily by the past and needs a better account of organizational change. Model III can become extremely complex and informationally demanding. The article therefore ends with a research agenda rather than a finished synthesis: foreign-policy analysis should make its conceptual models explicit, derive propositions more systematically, and combine strategic, organizational, and political variables in ways that improve both explanation and prediction.

\clearpage
\section{Jonathan Bendor and Thomas H. Hammond --- \emph{Rethinking Allison's Models}}

\subsection{Purpose and general critique}

Bendor and Hammond revisit Graham Allison's \emph{Essence of Decision} more than two decades after its publication. They acknowledge its enormous contribution: Allison showed how explicit models could generate alternative explanations of the same foreign-policy event and helped move the study of bureaucracy away from purely descriptive accounts. Their concern, however, is that the three famous models---rational actor, organizational process, and governmental politics---are less internally coherent than their influence suggests. The article therefore evaluates the models primarily as theories rather than reassessing the historical record of the Cuban Missile Crisis.

They develop five broad criticisms. First, the assumptions underlying Allison's models are often unclear. Second, many propositions are not logically derived from those assumptions. Third, some important propositions, especially in Model II, are logically wrong. Fourth, the models do not strike an appropriate balance between simplicity and complexity: Model I is too thin, whereas Model III is too thick. Fifth, Allison sometimes misrepresents the intellectual traditions from which the models are supposedly drawn. This matters because many readers have used \emph{Essence of Decision} as an introduction to rational choice, organization theory, and bureaucratic politics.

\subsection{A typology of policymaking models}

To clarify the problem, Bendor and Hammond construct a typology based on four dimensions: the number of decision makers, whether multiple actors have common or conflicting goals, whether actors are perfectly or imperfectly rational, and whether information is complete or incomplete. The typology demonstrates that Allison's models combine several dimensions at once, making it hard to know which assumption drives a particular prediction.

Model I appears to assume a single, perfectly rational actor, often with complete information. Models II and III are harder to classify. Model II is associated with bounded rationality and organizational routines, but Allison also introduces conflicting organizational goals. Model III clearly contains multiple actors with conflicting goals, yet it is ambiguous about whether those actors are fully rational or boundedly rational. It also incorporates misperception, incomplete information, bargaining, foul-ups, and incrementalism. As a result, Models II and III overlap analytically more than Allison's presentation implies. Before combining them, the authors argue, one should clarify exactly what each assumes.

\subsection{Model I: an unfairly weak version of rational choice}

Bendor and Hammond argue that Allison's Model I is not a fair test of rational-actor analysis because it is an unnecessarily primitive version of rational choice. Rationality does not require a single goal. An actor can pursue multiple, even competing objectives so long as preferences are consistently ordered. In the Soviet decision to deploy missiles in Cuba, for example, concern with the strategic nuclear balance could have coexisted with goals such as defending Cuba or reducing U.S. prestige. Showing that missile deployment was not optimal for every subsidiary goal does not prove those goals were irrelevant.

The model is also excessively static. Allison's rational actor chooses once at a single moment, whereas real foreign-policy decisions have consequences across time. A proper rational model can incorporate intertemporal trade-offs and discounting. Likewise, Model I largely neglects uncertainty. Rational-choice theory can analyze risk, imperfect knowledge of external events, and uncertainty about the preferences or capabilities of other states. Actors may be risk-averse, risk-neutral, or risk-seeking, and these attitudes affect choice.

Most importantly, Allison neglects strategic interaction. In international relations, the strongest justification for treating states as unitary rational actors is precisely that doing so allows the analyst to study strategic relations among states. Yet Allison's formal statement of Model I resembles individual decision theory, in which an isolated actor selects the best option from a fixed menu. Crisis bargaining is instead a game: each state's optimal action depends on the expected response of the other. Concepts such as Nash equilibrium, multiple equilibria, Pareto-inferior equilibria, and mixed strategies demonstrate that rationality does not imply a simple mapping from preferences to outcomes. A Prisoner's Dilemma can produce an outcome both sides dislike, while Chicken can contain multiple equilibria with very different distributions of gains. Strategic uncertainty can arise even when actors know the structure of the game.

The authors therefore reject the common lesson that the Cuban crisis proves the weakness of unitary rational-actor explanations. Allison's particular version of Model I was so underdeveloped that it should have been expected to perform poorly. A richer rational-choice model, especially one incorporating game theory and incomplete information, could generate much more powerful explanations.

\subsection{Model II: routines do not necessarily imply simple behavior}

Bendor and Hammond regard Allison's treatment of organizational routines as pioneering but criticize his interpretation of bounded rationality. Allison moves from the plausible claim that individuals use simple decision rules and SOPs to the stronger claim that organizational behavior will therefore be simple, rigid, incremental, and predictable. The authors argue that this inference is invalid. Complex behavior can emerge from simple rules, especially when many rules interact. Even a single adaptive decision maker using simple heuristics can produce complicated patterns, and an organization can combine many specialized routines into sophisticated behavior.

Allison also underestimates the positive functions of organizations. Simon and March did not view routines merely as constraints. Division of labor, specialization, redundancy, and organizational procedures can compensate for the cognitive limitations of individuals. Organizations can process more information and perform more complex tasks than any single member. SOPs can therefore \emph{enable} action as well as constrain it. Whether routines improve or degrade performance depends on the environment and on how well the routines match the problems faced.

The proposition that organizational behavior at $t+1$ will be only marginally different from behavior at $t$ is consequently too strong. Stable routines can generate nonlinear or surprising aggregate outputs, and organizations can recombine elementary programs in novel ways. The important theoretical question is not simply whether routines exist---all organizations require them---but how particular structures, procedures, information systems, and adaptation mechanisms shape performance.

\subsection{Model III: bargaining, hierarchy, and excessive complexity}

The authors' central criticism of Model III is that Allison treats bargaining as pervasive without explaining when superiors actually need to bargain with subordinates. The president possesses formal authority, appoints senior officials, and can in principle dismiss them. Why, then, must policy always be the result of bargaining among quasi-independent players? Bureaucratic-politics theory should identify the conditions under which hierarchy fails to produce command and bargaining becomes necessary.

Bendor and Hammond identify two especially important sources of subordinate influence. First, subordinates may possess political support outside the executive branch---in Congress, interest groups, professional communities, or other constituencies---which makes it costly for the president to override them. Second, subordinates may enjoy informational or technical advantages. They can influence which problems reach the president, which options are presented, and how decisions are implemented. Yet these advantages are variable. Presidential attention, alternative information channels, clear preferences, and control over appointments can reduce subordinate power. Bureaucratic politics should therefore be treated as conditional rather than universal.

Hierarchy itself is insufficiently theorized in Allison's Model III. Organizational structure determines who communicates with whom, which comparisons are made, which issues reach senior decision makers, and in what sequence alternatives are evaluated. Different hierarchies can therefore produce different policies even with identical actors and preferences. The interaction between formal authority and specialized expertise is a central feature of bureaucracy, but Allison does not develop it systematically.

The famous proposition that ``where you stand depends on where you sit'' also has ambiguous theoretical status. Position can influence preferences, but individuals are not reducible to their offices, and the proposition needs conditions and mechanisms if it is to be testable. More broadly, many Model III propositions are not clearly deduced from the model's premises.

Finally, Model III includes so many variables---positions, personalities, perceptions, bargaining advantages, organizational interests, deadlines, coalitions, personal commitments, and much more---that it risks becoming an ``analytical kitchen sink.'' A useful model must simplify. If a theory contains almost every potentially relevant feature of a case, it may describe events richly but cannot generate clear, falsifiable predictions. The authors suggest concentrating on a smaller number of structural variables, especially hierarchy, authority, expertise, informational asymmetry, and external political support.

\subsection{Conclusion}

Bendor and Hammond preserve Allison's methodological achievement while rejecting the theoretical authority often granted to his specific models. \emph{Essence of Decision} taught scholars to formulate explicit alternatives and to ask how different assumptions produce different explanations. That remains valuable. But the actual models require substantial reformulation.

Their final assessment is fivefold. Allison's models do not accurately represent all the literatures on which they are based; Model I does not give rational choice a fair test; the assumptions separating Models II and III are ambiguous; important propositions are not rigorously derived, and some are incorrect; and empirical support for a proposition teaches little about a model if the proposition is not logically connected to that model. Advances after 1971 make some weaknesses more visible, but many problems were present from the beginning. The article therefore calls for models that are explicit enough to reveal their assumptions, simple enough to yield deductions, rich enough to capture relevant strategic and organizational processes, and precise enough to be tested.

\clearpage
\section{Bruce Bueno de Mesquita and Alastair Smith --- \emph{Domestic Explanations of International Relations}}

\subsection{From unitary states to political leaders and institutions}

Bueno de Mesquita and Smith review the growth of research linking domestic politics to international relations. Their central claim is that foreign policy cannot be adequately understood by treating states as unitary actors with a single national interest. Leaders, advisers, supporters, opposition groups, and citizens may have different preferences, and political institutions determine whose interests matter for policy. The modern linkage literature therefore explains international behavior by examining leaders' incentives for political survival, principal-agent problems, domestic accountability, and institutional variation across regimes.

The political-economy approach assumes that leaders generally want to remain in office. They may promote national welfare when doing so supports that goal, but leader and citizen interests can diverge. The relevant domestic constituency varies across systems: in a democracy it may be a broad electorate, while in an autocracy it may consist of a small group of generals, party officials, clans, or cronies. Foreign policy must therefore be incentive-compatible with the political requirements leaders face at home. This perspective connects domestic institutions to war, crisis bargaining, nation building, foreign aid, and economic sanctions.

\subsection{Audience costs and international disputes}

A major strand of the literature examines \emph{audience costs}: domestic political penalties that leaders may suffer if they publicly escalate a dispute and then back down. Fearon's model argues that these costs are higher for democratic leaders because they are accountable to broad constituencies. Public threats can therefore ``tie the hands'' of democratic leaders and signal resolve more credibly. The implication is not simply that democracies are more belligerent. Rather, democratic leaders should be more selective about which crises they enter and which threats they make. If they do issue a threat, opponents may be more likely to believe it.

Later work endogenizes audience costs rather than assuming them. Smith links them to voters' uncertainty about leader competence and to reelection incentives. Slantchev shows that audience costs need not rise monotonically with democracy. Ashworth and Ramsay frame the problem in terms of moral hazard, asking how citizens optimally constrain leaders who may bluff to extract concessions. Schultz adds a domestic opposition party: because an opposition can criticize weak foreign-policy initiatives, its support or silence can itself convey information about the government's resolve. Across these models, the important point is that domestic institutions affect both which disputes are selected into observation and how signals are interpreted internationally.

\subsection{The democratic peace}

The most famous empirical regularity associated with domestic institutions is that democracies rarely fight wars against one another. The article reviews several explanations. Normative theories argue that democracies internalize habits of compromise and extend them to relations with other democracies. Institutional theories instead focus on incentives and selection effects.

Selectorate theory explains foreign policy through the size of the \emph{winning coalition}, the group whose support is necessary for a leader to remain in power, and the broader selectorate from which that coalition is drawn. Leaders dependent on large coalitions must provide relatively broad public benefits, including successful foreign policy. Defeat in war is therefore especially dangerous to democratic leaders. They become highly selective, fighting mainly when prospects of victory are strong, and once at war they have incentives to commit substantial resources to winning. Autocrats, whose survival may depend on a smaller coalition rewarded with private benefits, can tolerate policy failure more easily if they preserve resources for key supporters.

This logic helps explain why democracies are unattractive targets for other democracies. Both sides would need high confidence of victory, a condition that is difficult for two opponents to hold simultaneously. Democracies therefore tend to negotiate rather than fight one another, even though democracies are not inherently pacifist and may fight autocracies or weaker opponents. The review also discusses competing arguments involving issue salience, capitalism, leader experience, and the personal consequences of losing office. The larger lesson is that regime type shapes conflict through specific institutional incentives rather than through an undifferentiated ``national interest.''

\subsection{Diversionary conflict}

Domestic political weakness can also affect the timing and form of foreign policy. Diversionary-war arguments propose that leaders may initiate international confrontations to improve domestic standing. Research shows that election cycles, reelection prospects, term limits, and economic conditions can alter incentives for international action. Yet poor domestic conditions do not mechanically cause war. A politically vulnerable leader may instead make concessions in negotiations if an international agreement can generate political benefits at home, while foreign opponents may strategically accommodate a vulnerable but powerful state to avoid being targeted.

Models centered on leader competence clarify these mixed findings. International conflict can serve as a signal of competence only under certain conditions, for example when the opponent is strong enough that success is informative and when the incumbent values remaining in office sufficiently. The domestic-international linkage therefore produces conditional predictions rather than a simple claim that unpopular leaders start wars.

\subsection{Nation building and foreign aid}

The domestic perspective generates pessimistic expectations about externally imposed democratization. Democratic interveners are accountable to their own constituents and therefore seek foreign policy outcomes those constituents value. If the preferences of citizens in the target state conflict with the intervener's preferred policies, installing a fully accountable democracy can undermine the intervener's objectives. A dependent autocratic regime may be easier to influence because its leader needs to satisfy only a small coalition and can exchange policy concessions for resources used to reward supporters.

This logic also explains patterns of foreign aid. Aid involves donor leaders, donor constituents, recipient leaders, and recipient citizens. Donor governments provide aid when the foreign-policy concessions they can purchase are worth more politically than alternative domestic uses of the resources. Autocratic recipients are often attractive bargaining partners because they can trade policies for resources without needing broad domestic approval. Aid can therefore enhance the survival of autocrats and may do little for ordinary citizens. The literature reviewed by the authors finds strong connections between strategic interests, aid allocations, voting in international organizations, and access to IMF or World Bank resources. In this framework, aid is often better understood as a tool for buying policy compliance than as an instrument designed primarily for development.

\subsection{Economic sanctions}

The sanctions literature illustrates the importance of strategic selection. Simply asking how often imposed sanctions ``work'' is misleading. If a threatened sanction would be sufficiently painful, the target may concede before sanctions are imposed; the most effective sanctions therefore disappear from datasets of actual sanctions. Conversely, governments may impose weak sanctions that they know are unlikely to compel because imposing them yields domestic political benefits, such as appearing firm or satisfying interest groups.

Domestic institutions also affect the political consequences of sanctions. Leaders are not equally vulnerable to economic pain, because different regimes distribute costs differently and leaders depend on different constituencies. Sanctions may destabilize some leaders while allowing others to shift costs onto politically irrelevant groups. Their effectiveness therefore depends not only on aggregate economic damage but also on how that damage affects the coalition essential to the target leader's survival and on the sender's own domestic incentives.

\subsection{Conclusion}

The review concludes that the domestic-politics research program has substantially improved explanations of international behavior. Its key contribution is to replace the fiction of a homogeneous state with models of self-interested leaders operating under institutional constraints and interacting with supporters, opponents, and foreign actors. Audience costs, selectorate institutions, electoral incentives, political survival, and principal-agent problems help explain patterns that unitary-actor theories either miss or treat as exogenous. The authors therefore reject the idea that international relations constitute a separate realm of ``high politics.'' Foreign policy is made by political leaders, and domestic institutions systematically shape both their choices and the international consequences of those choices.

\clearpage
\section{Robert D. Putnam --- \emph{Diplomacy and Domestic Politics: The Logic of Two-Level Games}}

\subsection{Domestic-international entanglement}

Putnam's starting point is that domestic politics and international relations are mutually entangled. Asking whether domestic politics determines international relations or the international system determines domestic politics is therefore misguided; both causal directions operate under different conditions. The theoretical challenge is to explain when and how they interact. Putnam develops the metaphor of the \emph{two-level game} to show how national leaders must negotiate simultaneously with foreign counterparts and with domestic constituencies.

The 1978 Bonn economic summit provides the motivating example. The United States sought more expansionary policies from Germany and Japan, while Germany and Japan pressed the United States to address inflation, the dollar, and energy consumption. The summit produced a broad package: Germany accepted additional fiscal stimulus, Japan promised stronger growth and import expansion, and the United States committed to oil-price decontrol and other measures. Putnam argues that these changes cannot be explained either as pure responses to international pressure or as purely domestic choices. In each country, important domestic factions already favored policies demanded from abroad but lacked sufficient power to enact them alone. International pressure strengthened those factions, while domestic divisions made international concessions politically possible. The agreement succeeded because pressures at the two levels reinforced one another.

This example exposes the limits of both the ``Second Image,'' in which domestic politics causes international behavior, and the ``Second Image Reversed,'' in which international pressures cause domestic outcomes. Both are partial-equilibrium approaches. Putnam instead seeks a framework in which domestic and international bargaining are analyzed simultaneously.

\subsection{The two-level game}

At \emph{Level I}, national negotiators bargain with one another over an international agreement. At \emph{Level II}, domestic groups, parties, legislators, bureaucracies, and other constituents decide whether the agreement is politically acceptable and, where necessary, formally ratify it. The distinction is analytical rather than strictly chronological. Domestic bargaining shapes the mandate and constraints of international negotiators before and during Level I, while expectations about international bargaining affect domestic coalitions. Negotiations may move repeatedly between the two levels.

The chief negotiator sits at both tables. Internationally, the leader seeks an agreement with foreign counterparts. Domestically, the same leader must maintain a coalition capable of accepting or ratifying that agreement. A move that is attractive at one table may be costly at the other. A generous international concession may facilitate agreement abroad but provoke domestic defeat; a hard domestic position may preserve political support but make Level I agreement impossible. Skilled negotiators can sometimes use pressure at one level to change coalitions at the other.

Putnam distinguishes \emph{voluntary} from \emph{involuntary defection}. Voluntary defection occurs when an actor knowingly violates an agreement it could have honored. Involuntary defection occurs when a negotiator genuinely supports an agreement but cannot obtain domestic ratification or implementation. International cooperation therefore depends not only on whether governments reach bargains but also on whether those bargains fall within the domestic political limits of all participants.

\subsection{Win-sets and ratification}

The core concept is the \emph{win-set}. For a given country, the Level II win-set is the set of all possible Level I agreements that would obtain the necessary domestic approval. International agreement is possible only where the win-sets of the negotiating states overlap.

Win-set size has two crucial effects. First, larger win-sets make international agreement more likely because they increase the range of deals that can be ratified on both sides. Small win-sets raise the risk of involuntary defection and can produce bargaining breakdown if there is no overlap. Second, however, a small win-set can create international bargaining power. A negotiator may credibly say, ``I would like to concede, but my domestic constituents will not accept it.'' If the other side wants an agreement, it may have to move closer to the constrained negotiator's preferred outcome. This produces a paradox: domestic weakness can become international strength.

The advantage has limits. Shrinking a win-set may improve a state's bargaining position while overlap remains, but if the win-set becomes too small, the feasible bargaining range disappears and no agreement can be ratified. Putnam's diagram of two contracting win-sets illustrates this trade-off: moderate domestic constraints can improve distributional outcomes, while excessive constraints create deadlock.

Perceptions of win-sets are also strategically important. Negotiators may exaggerate domestic opposition in order to resist concessions, while foreign governments try to estimate which claims about domestic constraints are genuine. Uncertainty about the opposing win-set can generate tactical opportunities but also increase the risk of failed ratification.

\subsection{Determinants of win-set size: preferences and coalitions}

Putnam identifies three broad determinants of win-set size: Level II preferences and coalitions, Level II political institutions, and Level I negotiating strategies.

The first concerns the distribution of domestic preferences. If the costs of no agreement are low for important domestic groups, those groups will demand more from an international bargain and the win-set will be smaller. If no agreement would be very costly, more domestic actors will accept compromise and the win-set expands. The political meaning of ``no agreement'' therefore matters as much as the substantive terms of a proposed deal.

Domestic preference patterns can take different forms. On some issues, interests are relatively homogeneous and actors mainly disagree over how much to demand from abroad: there are domestic ``hawks'' and ``doves.'' On other issues, interests are heterogeneous, so different groups care about different components of a package. Heterogeneity can actually facilitate international cooperation because negotiators can construct packages that combine gains and losses across issues. A policy unacceptable as a single issue may become ratifiable when linked to another issue that benefits the opposing domestic coalition.

This possibility produces \emph{synergistic linkage}. International negotiations can alter domestic coalitions by bundling issues in new ways. Such linkage differs from simple bargaining tricks because the package changes the political feasibility of policies at Level II. The Bonn summit worked partly through this mechanism: international commitments supplied domestic factions with leverage and allowed governments to present a broader package of reciprocal concessions.

The relative political power of domestic groups also matters. Groups that are more organized, intense, or institutionally privileged may constrain ratification disproportionately. Because the composition of the active constituency differs across issues, a country's win-set is not a fixed national characteristic. It varies with the distribution of domestic costs and benefits associated with each negotiation.

\subsection{Domestic institutions and ratification rules}

The second determinant is political institutions. Formal ratification requirements can directly change win-set size. Requiring a simple majority generally creates a larger win-set than requiring a two-thirds majority or unanimity. Stronger ratification hurdles can therefore strengthen a negotiator's international bargaining position but also increase the chance that no agreement will survive Level II.

Informal institutions are equally important. Party discipline, coalition rules, legislative procedures, bureaucratic authority, and traditions of consensus shape what leaders can plausibly deliver. Strong party discipline can expand a leader's win-set by making legislative approval easier. Conversely, systems requiring broad consensus can narrow it.

Putnam uses this logic to reinterpret ``state strength.'' A state whose central executive is insulated from domestic pressure has a large win-set because the leader can ratify many agreements. That autonomy facilitates international cooperation, but it can weaken bargaining leverage: the executive cannot credibly claim that domestic actors prevent it from making concessions. A politically constrained negotiator may have less freedom at home but greater power at the international table. Institutional strength is therefore not unambiguously advantageous.

Real negotiations may involve more than two levels. An international agreement can require approval by a national legislature, parties within a coalition, administrative agencies, regional governments, or organized interests. In the European Community, for example, international decisions can pass through several nested levels of ratification. Each additional veto point changes the effective win-set and complicates bargaining.

\subsection{Level I strategies: changing the other side's win-set}

The third determinant concerns strategies used by international negotiators. A negotiator does not merely observe the foreign win-set; he or she may try to alter it. Side-payments, issue linkage, threats, symbolic gestures, and selective concessions can be targeted at politically pivotal groups in the other country. Successful diplomacy therefore requires knowledge of the opponent's domestic politics.

A concession with little value to the foreign government in the abstract may be crucial if it helps that government assemble a ratifying coalition. Conversely, a threat that imposes costs on politically marginal groups may have little effect, while a smaller but carefully targeted cost can transform the opponent's Level II balance. International bargaining power depends on the domestic incidence of external actions, not merely on their aggregate national effect.

This creates the possibility of transnational alliances. A foreign negotiator may effectively cooperate with domestic groups inside the other state that favor the same international agreement. Such alliances need not be formal. International pressure can supply arguments and resources to domestic minorities, while domestic actors may actively invite foreign pressure to strengthen their own position.

\subsection{Reverberation}

Putnam calls the effect of international pressure on domestic preferences and coalitions \emph{reverberation}. Foreign demands do not simply constrain a fixed domestic game; they can change it. At Bonn, external pressure strengthened domestic supporters of fiscal expansion in Germany and Japan and supporters of energy reform in the United States. International negotiation changed the domestic definition of what was politically feasible and, in some cases, what actors believed to be desirable.

Reverberation can be material or cognitive. Materially, international offers or threats can alter domestic costs and benefits. Cognitively, information and arguments from foreign governments or international institutions can persuade undecided actors, legitimize a policy, or strengthen a domestic minority. Positive reverberation expands the win-set and facilitates cooperation. Negative reverberation produces backlash and shrinks the win-set, especially when pressure comes from a distrusted adversary.

The concept is theoretically important because it prevents analysts from modeling domestic and international politics as independent sequential games. If international strategies can reshape domestic preferences, then Level II outcomes are not exogenous inputs into Level I bargaining. The two levels are genuinely interconnected.

\subsection{The role of the chief negotiator}

For analytical simplicity, the basic model initially treats the chief negotiator as an honest representative of domestic constituents. Putnam then relaxes this assumption. Leaders have their own preferences, political ambitions, reputations, and investments in domestic coalitions. The principal-agent relationship between leaders and constituents therefore matters.

A negotiator may value an international agreement because it enhances domestic political standing, may oppose an agreement that would disrupt an existing support coalition, or may hold personal policy preferences that differ from the median domestic actor. The leader also normally possesses an effective veto: even an agreement that could be ratified may never be submitted if the negotiator dislikes it. Woodrow Wilson's refusal to accept reservations to the League of Nations treaty illustrates how a leader can block an agreement that might otherwise lie within the domestic win-set.

Existing domestic coalitions create ``fixed investments.'' A deal that requires a leader to abandon core supporters and construct a new coalition can be politically costly even if a theoretical majority exists for ratification. Consequently, simple median-voter models may misrepresent the constraints facing actual negotiators. The identity and coalition of the leader matter.

The gap between leader and constituent preferences also creates additional international leverage. If foreign negotiators know that domestic constituents strongly favor an agreement while the leader is reluctant, they may exploit the leader's vulnerability. Conversely, a domestically weak leader facing hard-line constituents may gain bargaining power by credibly claiming that only a favorable agreement can survive politically.

\subsection{Conclusion}

Putnam's two-level framework integrates domestic and international politics without reducing one to the other. It recognizes that ``national interests'' are contested domestically and that national leaders must reconcile international bargains with domestic political survival and ratification. This produces strategic opportunities and dilemmas absent from unitary-actor models.

The framework highlights several recurring mechanisms: voluntary versus involuntary defection; the role of homogeneous and heterogeneous domestic interests; synergistic issue linkage; the paradox that domestic institutional weakness can strengthen international bargaining leverage; targeted threats, concessions, and side-payments; strategic uncertainty about win-sets; reverberation across levels; and conflicts between the chief negotiator's interests and those of domestic constituents.

The broader implication is that successful diplomacy is often not a matter of finding a compromise between two fixed national preferences. It involves constructing agreements that are simultaneously viable in several political arenas. International negotiators bargain not only with one another but also, indirectly, with each other's domestic coalitions. Because the size and composition of win-sets can be changed by institutions, strategies, information, and issue linkage, domestic politics is part of the bargaining process itself. Putnam therefore presents two-level games as a general research framework for explaining both international cooperation and bargaining failure and calls for systematic empirical work to develop and test its propositions.
